\documentclass[11pt,a4paper]{article}
\usepackage[utf8]{inputenc}
\usepackage[T1]{fontenc}
% Scommentare in locale se la texlive li include:
%\usepackage[italian]{babel}
%\usepackage{lmodern}
%\usepackage{microtype}
\usepackage[margin=2.6cm]{geometry}
\usepackage{booktabs}
\usepackage{amsmath}
\usepackage{amssymb}
\usepackage{array}
\usepackage{xcolor}
\usepackage{listings}
\usepackage{hyperref}
\usepackage{enumitem}
\hypersetup{colorlinks=true, linkcolor=blue!50!black, urlcolor=blue!50!black}

\lstset{
  basicstyle=\ttfamily\footnotesize, breaklines=true, frame=single,
  framerule=0.2pt, backgroundcolor=\color{gray!6}, showstringspaces=false
}

\newcommand{\code}[1]{\texttt{#1}}
\newcommand{\flow}[1]{\texttt{\small #1}}

\title{\textbf{Pipeline SDMX--ISTAT per popolazioni sintetiche comunali e sub-comunali}\\[2mm]
\large Note tecniche --- v0.4}
\author{M.~Degli Esposti}
\date{19 luglio 2026 \quad (v0.1--v0.3: 18--19 luglio 2026)}

\begin{document}
\maketitle

\begin{abstract}
\noindent
Pipeline riproducibile dai dati pubblici ISTAT alla popolazione sintetica
MaxEnt con quartieri: Brescia K7C, $198\,259$ individui $\times$ 11 attributi
(\code{popolazione\_K7C\_naz.csv}). In v0.4: (xi) nazionalit\`a two-stage
gerarchica --- area UE/extra-UE condizionata a (zona, sesso) dalle sezioni,
paese condizionato a (area, sesso) dal comunale, allocazione esatta
largest-remainder; doppia contabilit\`a comunale$\leftrightarrow$sezioni
verificata a $\sim3$\textperthousand; mappa quartiere$\times$paese-top
qualitativamente corretta (Fiumicello/Chiesanuova a trazione Pakistan,
Primo Maggio/Don Bosco a trazione Romania); (xii) ricettario di esecuzione
end-to-end della pipeline (\S\ref{sec:recipe}). Risultati precedenti:
identit\`a esatte della base riconciliata (5$\times$); regole di feasibility
(completamento a partizione; consistenza di tutti i margini sovrapposti + IPF
come armonizzatore); formula del pavimento stocastico PCD replicata 3 volte;
stimatore-pool vs distribuzione appresa.
\end{abstract}

\tableofcontents

%=======================================================================
\section{Scopo e contesto}
Dato un codice comune ISTAT, produrre il miglior constraint set dai dati
pubblici e la popolazione MaxEnt corrispondente; caso guida Brescia (SIN
Caffaro) con zona a 33 quartieri. Esecuzione locale (WSL2, env \code{ml});
solver dal repo \code{maxent-popsynth-pcd}.

%=======================================================================
\section{Endpoint SDMX \code{esploradati}: comportamenti non documentati}
\label{sec:quirks}
\begin{enumerate}[leftmargin=*,itemsep=1pt]
  \item limite $\sim$260 char/segmento URL (http.sys) $\Rightarrow$ wildcard nel KEY;
  \item \code{startPeriod}/\code{endPeriod} ignorati $\Rightarrow$ filtro client-side
  (serie storiche complete gratis);
  \item SDMX-CSV via \code{Accept: application/vnd.sdmx.data+csv;version=1.0.0};
  \item rate limit $\sim$5/min: throttling 13\,s, retry con backoff;
  \item stime censuarie comunali \emph{non arrotondate} (decimali, additive).
\end{enumerate}

%=======================================================================
\section{Le fonti: tre laghi pi\`u uno}
\label{sec:laghi}
Anagrafe (\flow{DCIS\_*}, singolo anno d'et\`a); censimento permanente
comunale (\flow{DF\_DCSS\_*}, ipercubo unico, granularit\`a nel suffisso);
indagini campionarie (solo regionali $\to$ condizionali two-stage); canale
sub-comunale:
\begin{itemize}[leftmargin=*,itemsep=1pt]
  \item ASC aggregate via SDMX (\flow{DF\_BULK\_SUBCOM*}): solo capoluoghi
  Citt\`a Metropolitane;
  \item \textbf{sezioni di censimento}, tutti i comuni: bulk
  \code{Dati\_regionali\_<anno>.zip} (xlsx per regione + tracciato) su
  \code{esploradati.istat.it/databrowser/DWL/PERMPOP/SUBCOM/}; edizioni 2021
  e 2023 (rilascio annuale a regime);
  \item scoperta chiave: \code{COM\_ASC1/2/3} nel file sezioni --- i codici
  delle aree sub-comunali sono gi\`a assegnati sezione per sezione per tutti
  i comuni partizionati (Brescia: 33 valori $=$ quartieri; spatial join
  superfluo).
\end{itemize}
Variabili del tracciato: sesso$\times$et\`a quinquennale, istruzione a 5
livelli per sesso (9+), occupati 15--64 per sesso, ITA/stranieri per
sesso$\times$macro-et\`a, \textbf{UE/extra-UE per sesso} (ST16--21),
background migratorio, famiglie, abitazioni.

%=======================================================================
\section{Brescia: sezioni, quartieri, audit}
1\,822 sezioni, 33 quartieri; sezione fittizia 888888x (549 senza fissa
dimora) tenuta nel quartiere assegnato \code{17029004} (opzione a).
\textbf{Quinta identit\`a esatta}: sezioni $\equiv$ censimento comunale 2023
$\equiv$ anagrafe 1/1/2024 $= 198\,259$ (stranieri $37\,478$); additivit\`a
interna piena. Sezioni in interi (RBI$\leftrightarrow$RSBL), comunale in
stime decimali: scarti di unit\`a assorbiti dal disegno condizionale.
Ancoraggio K7C: \code{--anno 2024}. Sanity: top-5 quota stranieri $=$
Fiumicello (0.371), Centro Storico Nord, Primo Maggio, Don Bosco,
Chiesanuova --- l'orbita SIN Caffaro.

%=======================================================================
\section{K7C: blocchi di zona, feasibility, IPF}
\label{sec:k7c}
\code{zona} in testa a \code{VAR\_ORDER}; $|X| = 33\times5\,376 = 177\,408$;
ariet\`a piena dove il dato la sostiene (Z2:
zona$\times$sesso$\times$et\`a$\times$cittadinanza). Principio: ogni tabella
di zona entra come $P(\text{zona}\mid\text{gruppo})\times$conteggi comunali.
Assunzioni (2) quota di zona costante entro quinquennale/macro-classe;
(3) $P(\text{zona}\mid\text{sesso},\text{occ})$ costante sui bin 15--64.

\medskip\noindent\textbf{Regola di consistenza generalizzata} (dal collasso
discesa$\to$risalita del primo fit): con vincoli di uguaglianza devono
coincidere numericamente \emph{tutti} i margini impliciti sovrapposti tra
blocchi, non solo quelli verso la spine. \textbf{Fix: IPF come armonizzatore
di tabelle} (righe da Z1, colonne dal blocco comunale, per gruppo); Z4 con
cap su Z1 e ridistribuzione. Audit post: Z-vs-Z $\sim10^{-11}$,
Z-vs-comunale $=0$.

%=======================================================================
\section{Fit K7C e caratterizzazione GibbsPCD}
\begin{center}
\small
\begin{tabular}{@{}lrrrrr@{}}
\toprule
 & \textbf{MRE}$_{\alpha>0}$ & \textbf{H (nat)} & \textbf{supporto} & \textbf{KL vs exact} & \textbf{tempo}\\
\midrule
Esatto (L-BFGS) & $3.6\cdot10^{-4}$ & 9.157 & 130\,161/177\,408 & --- & 351\,s\\
GibbsPCD ($N{=}2\cdot10^5$, numba) & $3.1\cdot10^{-2}$ & 9.592 & 171\,864 & $1.2\cdot10^{-1}$ & 1\,070\,s\\
\bottomrule
\end{tabular}
\end{center}
Esclusioni ($\alpha=0$) fuori dai solver, post-hoc sul supporto (5\,544
celle). \textbf{Pavimento stocastico replicato} (terzo punto):
predetto $0.0650$, osservato $0.0634$ ($N=2\cdot10^5$); formula
$\text{MRE}_{\min}\approx m^{-1}\sum_j\sqrt{(1-\alpha_j)/(\alpha_j N)}$
calcolabile dal solo CS. \textbf{Stimatore-pool vs distribuzione appresa}:
MRE su frequenze del pool $0.063$ (al pavimento), MRE su $p_\lambda$
$0.031$: i $\lambda$ mediati da Adam battono il pavimento del singolo
snapshot di $\sim2\times$. Decisione di produzione: esatto anche a K7C
($100\times$ pi\`u preciso, $1/3$ del tempo); GibbsPCD per domini non
enumerabili, $N$ dalla formula.

%=======================================================================
\section{Nazionalit\`a two-stage (livello 1)}
\label{sec:naz}
La nazionalit\`a (143 paesi) resta fuori dal MaxEnt (nessun incrocio
vincolabile; $|X|$ esploderebbe) ed entra come assegnazione gerarchica sul
campione (\code{assign\_nationality.py}):
\begin{enumerate}[leftmargin=*,itemsep=1pt]
  \item[(2a)] area $\in$ \{UE, EXTRA\_UE\} $\sim P(\text{area}\mid
  \text{zona},\text{sesso})$ dalle sezioni (ST17/18/20/21; UE$+$extraUE
  $=$ ST1 esatto, apolidi assorbiti in extra-UE);
  \item[(2b)] paese $\sim P(\text{paese}\mid\text{area},\text{sesso})$ dal
  comunale (\flow{DF\_DCSS\_POP\_DEMCITMIG\_TV\_3}, 2023, GENDER
  conservato); classificazione UE label-driven (23 paesi UE presenti a
  Brescia; Cechia/Lussemburgo/Malta assenti dal dato, non errori di match).
\end{enumerate}
Allocazione \textbf{esatta} (largest remainder) a entrambi i livelli, shuffle
seedato entro gruppo. \textbf{Assunzione (4)}: paese $\perp$ zona $\mid$
(area, sesso) --- la struttura territoriale entra al livello area; il limite
\`e visibile (dentro l'extra-UE stesso mix in tutti i quartieri).

\medskip\noindent\textbf{Validazione}: doppia contabilit\`a
comunale$\leftrightarrow$sezioni per area$\times$sesso coincide a
$\sim3$\textperthousand\ (canali indipendenti); top-10 campione $=$
Romania, Pakistan, Ucraina, India, Egitto, Cina, Albania, Moldova, Sri
Lanka, Filippine (ordine reale, Ucraina con asimmetria F via
condizionamento per sesso); paese-top per quartiere: Fiumicello e
Chiesanuova $\to$ Pakistan, Primo Maggio e Don Bosco $\to$ Romania ---
differenziazione emergente da $P(\text{area}\mid\text{zona},\text{sesso})$.

%=======================================================================
\section{Sintesi dei livelli e deliverable}
\begin{center}
\small
\begin{tabular}{@{}lrrrrr@{}}
\toprule
 & \textbf{anno} & $|X|$ & $m$ & \textbf{MRE esatto} & \textbf{H (nat)}\\
\midrule
K6C & 1/1/2025 & 5\,376 & 269 & $5.0\cdot10^{-4}$ & 5.842\\
K7C & 1/1/2024 & 177\,408 & 2\,513 & $3.6\cdot10^{-4}$ & 9.157\\
\bottomrule
\end{tabular}
\end{center}
\noindent\textbf{Deliverable}: \code{popolazione\_K7C\_naz.csv} ---
$198\,259$ individui $\times$ 11 attributi (zona, quartiere, sesso, eta,
stato\_civile, cittadinanza, istruzione, condizione, area, paese; ITL $\to$
paese Italia), ancorati all'1/1/2024, con 4 assunzioni dichiarate e
numerate.

%=======================================================================
\section{Ricettario: esecuzione della pipeline}
\label{sec:recipe}
Sequenza completa da zero a \code{popolazione\_K7C\_naz.csv}. Prerequisiti:
env \code{ml} attivo; repo \code{maxent-popsynth-pcd} clonato in
\code{\textasciitilde/progetti/}; \code{pip install lxml openpyxl}. Script
in \code{\textasciitilde/progetti/gsp/scripts/}. Le fasi 0--4 sono
\textbf{generiche per qualunque comune} (parametro \code{<codice>},
ITTER107); le fasi 5--8 sono generiche nel disegno ma gli script attuali
hanno tre punti parametrizzati su Brescia, elencati in
\S\ref{sec:adapt}.

\medskip\noindent\textbf{Fase 0 --- catalogo (una tantum, riusabile).}
\begin{lstlisting}
python istat_catalog.py                 # scarica e indicizza 4884 dataflow
python istat_catalog.py istruzione      # esempio di grep esplorativa
\end{lstlisting}

\noindent\textbf{Fase 1 --- tavole comunali} (esplorazione staged, poi fetch;
$\sim$3 min per il throttling).
\begin{lstlisting}
python fetch_comune.py <codice> --explore    # strutture: dimensioni, territorio
python fetch_comune.py <codice>              # 7 tavole -> data/comuni/<codice>/
\end{lstlisting}

\noindent\textbf{Fase 2 --- tavole-vincolo comunali} (per ogni ancoraggio che
serve; K7C usa \code{--anno 2024}).
\begin{lstlisting}
python build_constraints.py <codice> --anno 2024
python build_constraints.py <codice> --anno 2025    # opzionale, K6C recente
\end{lstlisting}
Verificare nel log: pop totale plausibile; report.md con MAE del raccordo
(atteso $\approx0$ post-2018).

\medskip\noindent\textbf{Fase 3 --- ConstraintSet K6C e fit} (opzionale ma
consigliato come regressione).
\begin{lstlisting}
python cs_build.py <codice> --anno 2025
python fit_cs.py <codice> --anno 2025 --eps 1e-8 --min-alpha 2e-4 --no-gibbs
\end{lstlisting}

\noindent\textbf{Fase 4 --- sezioni di censimento} (bulk, nessun rate limit).
\begin{lstlisting}
mkdir -p ~/progetti/gsp/data/submun && cd ~/progetti/gsp/data/submun
wget https://esploradati.istat.it/databrowser/DWL/PERMPOP/SUBCOM/Dati_regionali_2023.zip
unzip Dati_regionali_2023.zip "Dati_regionali_2023/R<NN>_<Regione>_2023_sezioni.xlsx" \
      "Dati_regionali_2023/TRACCIATO FILE REGIONALI.xlsx"
\end{lstlisting}
Estrarre il comune dal file regionale (filtro \code{PROCOM}, xlsx $\to$
csv) e verificare: n.\ sezioni, \code{COM\_ASC1} popolata e con quante zone,
sezioni fittizie 888888x, P1 vs popolazione attesa.

\medskip\noindent\textbf{Fase 5 --- tabelle di quartiere.}
\begin{lstlisting}
python build_zona_tables.py     # Z1-Z4 + zona_nomi -> data/comuni/<codice>/zona_2023/
\end{lstlisting}
Audit stampati: totali $=$ comunale, quota stranieri per quartiere (sanity
qualitativa).

\medskip\noindent\textbf{Fase 6 --- ConstraintSet K7C e fit.}
\begin{lstlisting}
python cs_build.py <codice> --anno 2024 --livello K7C
python fit_cs.py <codice> --anno 2024 --livello K7C --eps 1e-8 \
                --min-alpha 2e-4 --no-gibbs
\end{lstlisting}
Nel build verificare: somme blocchi (Z1--Z3 $=1.0000$), audit Z-vs-Z
$\sim10^{-11}$ e Z-vs-comunale $=0$. Nel fit: MRE$_{\alpha>0}$ attesa
$\sim10^{-4}$--$10^{-3}$; se compare discesa$\to$risalita, c'\`e un margine
sovrapposto non armonizzato (\S\ref{sec:k7c}). Output:
\code{popolazione\_K7C.csv}.
(Facoltativo, caratterizzazione GibbsPCD: stesso comando senza
\code{--no-gibbs}, con \code{--pool} dimensionato dalla formula del
pavimento e \code{--numba}.)

\medskip\noindent\textbf{Fase 7 --- nazionalit\`a two-stage.}
\begin{lstlisting}
python assign_nationality.py <codice> --anno 2024
\end{lstlisting}
Verificare: classificazione UE stampata; doppia contabilit\`a
comunale$\leftrightarrow$sezioni; top paesi e paese-top per quartiere.
Output finale: \code{popolazione\_K7C\_naz.csv}.

\subsection{Adattamento a un comune diverso da Brescia}
\label{sec:adapt}
Tre punti oggi parametrizzati su Brescia, da generalizzare (candidati per una
v3 degli script):
\begin{enumerate}[leftmargin=*,itemsep=1pt]
  \item \textbf{estrazione sezioni}: il filtro \code{PROCOM} dal file
  regionale \`e un one-liner manuale; \code{build\_zona\_tables.py} e
  \code{assign\_nationality.py} leggono il path hardcoded
  \code{brescia\_sezioni\_2023.csv};
  \item \textbf{mappa nomi ASC} (\code{ASC\_NOMI}): fornita a mano per
  Brescia; per altri comuni va reperita dal comune stesso (i codici
  \code{COM\_ASC1} sono nel dato, i nomi no) e l'\code{assert len==33}
  parametrizzato;
  \item \textbf{presenza delle ASC}: verificare che \code{COM\_ASC1} sia
  popolata per il comune scelto; se il comune non ha partizioni
  amministrative dichiarate, la zona va costruita per aggregazione
  spaziale delle sezioni (basi territoriali 2021 + confini scelti).
\end{enumerate}
Con questi tre punti resi parametrici, la pipeline \`e end-to-end generica:
\code{<codice>} $\to$ \code{popolazione\_K7C\_naz.csv}.

%=======================================================================
\section{Da fare / prossimi passi}
\begin{itemize}[leftmargin=*,itemsep=1pt]
  \item[$\boxtimes$] Nazionalit\`a two-stage livello 1 --- v0.4.
  \item[$\square$] Generalizzazione multi-comune degli script sub-comunali
  (\S\ref{sec:adapt}) $+$ run di regressione su Torino 001272.
  \item[$\square$] Run B (pool $=$ popolazione, $\lambda$ caldi): misura
  diretta del pavimento stazionario; richiede warm start (subclass del
  solver) --- da ridiscutere prima.
  \item[$\square$] Via/civico entro quartiere (macchinario geografico v3).
  \item[$\square$] Seed esplicito nel kernel Numba.
  \item[$\square$] Variante Caffaro: esclusione sezione fittizia (549) con
  riscalatura.
  \item[$\square$] $\sigma$ campionarie del censimento; UE/extra-UE come
  raffinamento della cittadinanza nel CS; famiglie e background migratorio.
  \item[$\square$] Aggiornamento annuale: sub-comunale 2024 al rilascio.
\end{itemize}

\end{document}
